\begin{abstract}
Large language model (LLM) agents act through tool calls, and an agent that reads attacker-controlled content can be steered into calling tools against its user. We present \SystemName, a policy decision point that judges every proposed tool call before it runs and answers \verdict{allow}, \verdict{modify}, \verdict{ask} or \verdict{deny}. It fails closed on every error, timeout and malformed input. It combines canonicalisation, action tiers with argument conditions, typed parsing of network arguments, session budgets, a predictive check of containment effects against an asset inventory, session information flow (the ``Rule of Two''), and calibrated deferral of approval requests with a Learn-then-Test guarantee. We replay AgentDojo's ground-truth tool calls (\AdBenign{} benign tasks and \AdAttacks{} fully hijacked-agent attacks) through policies generated only from tool descriptions. Action tiers alone stop all but \AdTiersGoalUnreviewed{} of attack goals from running without a person, at \AdTiersAsks{} approval requests per benign task (\AdAskReduction{} fewer than asking for everything); conservative information-flow labels close the rest at \AdConsAsks{} requests. Calibration is not containment: if deferral auto-approved every request it is allowed to, \DefTiersGoalUnreviewed{} of attack goals would run unreviewed with tiers alone, and \DefConsGoalUnreviewed{} with information-flow labels, which keep exposure-driven sends with a person. We also find that agent hosts let tool calls through when a hook fails or times out, so enforcement must be written for each host's failure semantics, and that typed argument parsing removes failure modes of string rules, including a URL parser differential that string rules sent to a reviewer as in scope. We release the filter, policies, replay harness and per-call results.
\end{abstract}

\begin{IEEEkeywords}
LLM agents, prompt injection, tool use, safety filter, information flow control, conformal risk control, reference monitor
\end{IEEEkeywords}
